\section{Conclusion}

No single observational study resolves a causal question definitively.
The field builds general causal knowledge through \emph{triangulation}: accumulating credible local estimates across complementary designs, data sources, and populations~\cite{shadish2002experimental}---just as economics built its evidence base on minimum wage effects and returns to education through dozens of independent, credible studies that converged~\cite{angrist2010credibility}.
SE's diversity of ecosystems---different languages, platforms, team structures, development cultures---makes it well suited to this strategy: The same causal question can be studied in many independent settings, each with its own identification strategy, and convergence across them is far more persuasive than any single study claiming a universal effect.

The causal inference toolkit's greatest contribution to SE is not a specific method but a way of thinking: making identifying assumptions \emph{transparent and open to scrutiny} so that the strength of evidence is clearly communicated and the field can iteratively develop stronger designs.
We envision this adoption bringing productive advancements on a broad range of important intervention-outcome questions---from the effect of AI coding tools to the impact of development practices, governance structures, and ecosystem policies---moving the field from the blanket disclaimer that ``correlation does not imply causation'' toward the constructive discipline of asking: \emph{Under what assumptions does this evidence support a causal interpretation, and how can the next study relax those assumptions?}
